\documentclass[11pt]{article}
\usepackage[a4paper,margin=0.9in]{geometry}
\usepackage{amsmath,amssymb,amsthm}
\usepackage{booktabs}
\usepackage{hyperref}
\usepackage{microtype}
\usepackage[T1]{fontenc}
\usepackage{lmodern}
\usepackage{xcolor}

\newtheorem{theorem}{Theorem}
\newtheorem{proposition}[theorem]{Proposition}
\newtheorem{lemma}[theorem]{Lemma}

\title{Exact constrained covering arrays for the public \texttt{INDUSTRIAL\_8} benchmark}
\author{Unsolved Labs}
\date{Research Release R003 -- August 2026}

\begin{document}
\maketitle

\begin{abstract}
For the public CT Competition 2023 model \texttt{INDUSTRIAL\_8}, we determine the exact minimum number $N_t$ of valid configurations required to cover every valid $t$-way interaction for strengths $t=2,\ldots,6$:
\[
N_2=18,\qquad N_3=45,\qquad N_4=72,\qquad N_5=N_6=144.
\]
For strength $3$, the pinned 2023 competition results contain a best valid public suite of 54 tests. The released 45-row suite reduces this comparison by 9 tests (16.7\%) and is optimal for the frozen model. The lower bounds follow from a structural decomposition into nine disjoint mandatory classes; the matching upper bounds are explicit test suites. Two independent exact implementations replay every released witness against all valid interactions, and the nontrivial local five-row covering-array bound used at strength 3 is additionally formalized in Lean 4. The research artifact was generated with frontier AI; search and generation are outside the final correctness oracle.
\end{abstract}

\section{Introduction}
Combinatorial interaction testing seeks small test suites in which every relevant interaction among parameter values is exercised. In constrained models, only globally valid configurations and interactions that can occur in at least one valid configuration are required. The second CT Competition compared state-of-the-art generators on a public benchmark set across strengths 2 through 6 \cite{ctcompetition2023}. The public benchmark generator and its validation methodology are described by Bombarda and Gargantini \cite{bombarda2024}.

This paper studies one frozen evaluation model, \texttt{INDUSTRIAL\_8} \cite{citbenchmark2023}. The result is exact rather than heuristic: the repository includes explicit constructions and matching lower bounds for every competition strength. At strength 3 the pinned competition table records valid suites of sizes 55, 68, 68, and 54; entries of sizes 71 and 49 are marked invalid. Thus 54 is the best valid value in that frozen table, while the construction proved here has size 45.

The proof is deliberately separated from the search process that discovered the constructions. The final trust boundary consists of the frozen benchmark, the analytic lower bounds below, explicit CSV witnesses, two independent exact replay implementations, and a Lean proof of the most delicate local strength-3 covering-array fact.

\section{Frozen model and coverage convention}
The model has Boolean parameters $p_1,\ldots,p_{13}$ and a three-valued parameter $p_{14}\in\{v1,v2,v3\}$. The exact source file is frozen in the public repository. \texttt{SOURCE\_AUDIT.md} pins the upstream commit, Git blob, and SHA-256 digest:
\begin{quote}\small
commit: \path{7ed44c18a771b9bbdc444a3c9178a7a68510a53f}\\
Git blob: \path{4ede6ee93f6eea06927f972cf5ba22621ecc3013}\\
SHA-256: \path{01785c43dbce4381d928c064c14dfe3c3734237db8fa7c2e84d8647f8ed299b4}
\end{quote}

A \emph{valid $t$-way interaction} is a choice of $t$ distinct parameter positions together with values that occur jointly in at least one globally valid model configuration. A test suite has strength $t$ if it covers every valid $t$-way interaction. Let $N_t$ be the minimum cardinality of such a suite.

\section{Structural reduction}
The model constraints imply the following facts.

\begin{lemma}[Model structure]\label{lem:structure}
Every valid configuration satisfies:
\begin{enumerate}
  \item $p_{14}\neq v3$;
  \item $p_{14}=v1$ if and only if $p_3=\mathrm{true}$, and otherwise $p_{14}=v2$;
  \item among $p_5,\ldots,p_{12}$, at most one parameter is false;
  \item if $p_{13}=\mathrm{false}$, then $p_5,\ldots,p_{12}$ are all true;
  \item the assignment $p_1=\cdots=p_{13}=\mathrm{true}$ is forbidden.
\end{enumerate}
Consequently the model has exactly 159 valid configurations.
\end{lemma}

\begin{proof}
The first two statements are direct rewrites of the three constraints on $p_{14}$. The pairwise clauses among $p_5,\ldots,p_{12}$ forbid two simultaneous false values. The clauses $p_{13}=\mathrm{true}\lor p_i=\mathrm{true}$, $5\le i\le 12$, give the fourth statement. The final long disjunction gives the fifth. Counting then yields eight classes of 16 configurations with exactly one of $p_5,\ldots,p_{12}$ false, one class of 16 configurations with $p_{13}$ false, and 15 residual all-true middle configurations, for $8\cdot 16+16+15=159$.
\end{proof}

Call the first nine 16-element groups the \emph{mandatory classes}. They are disjoint, and within each one the four coordinates $p_1,\ldots,p_4$ range freely over $\{0,1\}^4$.

\paragraph{Class additivity.}
Any interaction that fixes the defining false parameter of a mandatory class can only be covered by a row from that class. Hence lower bounds established independently inside the nine mandatory classes add.

\section{Lower bounds}
Fix one mandatory class. Since $p_1,\ldots,p_4$ are free inside the class, interactions involving the class-defining false parameter force standard binary covering requirements on those four coordinates.

\subsection{Strength 2}
\begin{lemma}[Local one-way minimum]\label{lem:one}
At least two rows are required in each mandatory class for global strength 2.
\end{lemma}
\begin{proof}
Fixing the class-defining false parameter and one of $p_1,\ldots,p_4$ requires both Boolean values, so one row cannot suffice.
\end{proof}
Across the nine disjoint mandatory classes this gives
\[
N_2\ge 9\cdot 2=18.
\]

\subsection{Strength 3}
Fixing the class-defining false parameter and two of $p_1,\ldots,p_4$ requires a binary pairwise covering array on four columns.

\begin{lemma}\label{lem:ca5}
Four rows cannot pairwise cover four binary columns, while five rows suffice.
\end{lemma}

\begin{proof}
If four rows covered all four value pairs on every column pair, then for every pair of columns the pairs $00,01,10,11$ would each occur exactly once. After permuting rows and complementing columns, normalize two columns to $0011$ and $0101$. Up to complement, the unique balanced column orthogonal to both is $0110$, leaving no fourth balanced column orthogonal to all three. Thus four rows are impossible.

Five rows suffice via
\[
B=\{0000,0111,1011,1101,1110\}.
\]
The impossibility and the explicit five-row witness are also checked by Lean in \path{R003Formal/Industrial8.lean}.
\end{proof}

Thus each mandatory class needs at least five rows and
\[
N_3\ge 9\cdot 5=45.
\]

\subsection{Strength 4}
\begin{lemma}[Local three-way minimum]\label{lem:three}
At least eight rows are required in each mandatory class for global strength 4.
\end{lemma}
\begin{proof}
Fixing the class-defining false parameter and any three of $p_1,\ldots,p_4$ requires all $2^3=8$ triples.
\end{proof}
Thus
\[
N_4\ge 9\cdot 8=72.
\]

\subsection{Strengths 5 and 6}
\begin{lemma}[Local four-way minimum]\label{lem:four}
At least sixteen rows are required in each mandatory class for global strengths 5 and 6.
\end{lemma}
\begin{proof}
At strength 5, fixing the class-defining false parameter and all four free Boolean coordinates requires all $2^4=16$ assignments. At strength 6, add any parameter forced true in the class; the same 16 four-bit assignments remain necessary.
\end{proof}
Therefore
\[
N_5\ge 144,\qquad N_6\ge144.
\]

\begin{theorem}[Global lower bounds]\label{thm:lower}
For the frozen benchmark,
\[
N_2\ge18,\quad N_3\ge45,\quad N_4\ge72,\quad N_5\ge144,\quad N_6\ge144.
\]
\end{theorem}
\begin{proof}
Each bound is the corresponding exact class-local requirement multiplied by the nine disjoint mandatory classes.
\end{proof}

\section{Matching constructions}
All constructions use only mandatory-class rows.

For strength 2, use two complementary four-bit words per class and cycle four cut types across the nine classes. For strength 3, use translations of the five-row block $B$ by $0000$, $0001$, $0010$, and $0011$, cycling these translations across the classes. For strength 4, alternate the even- and odd-parity eight-word sets; each realizes every triple on four coordinates. For strengths 5 and 6, use all 16 four-bit words in every mandatory class.

The repository releases a concrete CSV witness for every strength. Exact replay gives:

\begin{center}
\begin{tabular}{rrrr}
\toprule
Strength & valid interactions & released rows & lower bound \\
\midrule
2 & 326 & 18 & 18 \\
3 & 2,168 & 45 & 45 \\
4 & 9,374 & 72 & 72 \\
5 & 28,192 & 144 & 144 \\
6 & 61,272 & 144 & 144 \\
\bottomrule
\end{tabular}
\end{center}

The lower and upper bounds therefore agree at every strength.

\begin{theorem}[Exact optima]\label{thm:main}
For the frozen CT Competition 2023 \texttt{INDUSTRIAL\_8} model,
\[
N_2=18,\qquad N_3=45,\qquad N_4=72,\qquad N_5=N_6=144.
\]
\end{theorem}
\begin{proof}
Theorem~\ref{thm:lower} gives the lower bounds. The released CSV witnesses have exactly those sizes, and both exact replay implementations verify complete coverage of every feasible interaction.
\end{proof}

\section{Verification and reproducibility}
The public repository contains three complementary verification layers.

\paragraph{Python exact verifier.}
\texttt{verify.py} checks the byte identity of the benchmark against the pinned upstream Git blob, reads the released witness CSVs, enumerates all 24,576 raw assignments, reconstructs the 159 valid configurations and all valid interactions, and verifies each witness against the matching analytic lower bound. It also checks the frozen 2023 comparison excerpt and the deterministic machine-readable verification report.

\paragraph{Independent C++ replay.}
\texttt{verify\_independent.cpp} does not call or import the Python verifier. It independently parses the parameter domains and Boolean constraint expressions from the frozen model text, evaluates every raw assignment, reads the same released witness CSVs, rebuilds every feasible interaction set, and separately brute-forces the class-local covering minima.

\paragraph{Lean.}
The Lean project formalizes Lemma~\ref{lem:ca5}: the impossibility of a four-row pairwise cover on four binary columns and correctness of the explicit five-row block. The impossibility theorem follows the structural bijection argument above; the explicit five-row witness is checked by ordinary finite \texttt{decide}. This is intentionally described as a partial formalization: the full model reduction and the other exact minima remain supported by the human argument and two exact finite replay paths. The declarations contain no \texttt{sorry} or \texttt{admit}; the public axiom audit is recorded in \texttt{formalization.yaml}.

All finite release checks are invoked by \texttt{./verify\_release.sh} and mirrored in GitHub Actions. The Lean project is pinned to Lean 4 / Mathlib v4.32.0 and is built with \texttt{lake build}; CI additionally invokes an independent Lean type checker with proof placeholders disallowed.

\section{Related work, comparison boundary, and limitations}
The benchmark family and generator are public research artifacts described in the software-testing literature \cite{bombarda2024}. The frozen benchmark is pinned to commit \texttt{7ed44c18a771} of the public benchmark-generator repository. The strength-3 numerical comparison is limited to the frozen 2023 results file at commit \texttt{3c8cdc1ac9eb} of the public competition repository. \texttt{SOURCE\_AUDIT.md} records the complete repository paths, commit identifiers, and Git blobs. The release does not claim that 54 remained the best result in every later source, nor does it claim a general theorem for arbitrary constrained industrial models. External specialist review and peer review remain pending.

\section*{AI-origin disclosure}
This research release was generated with frontier AI. Public correctness claims rely on the explicit proof and verification artifacts described above rather than on the generation process itself.

\begin{thebibliography}{9}

\bibitem{bombarda2024}
Andrea Bombarda and Angelo Gargantini.
\newblock Design, implementation, and validation of a benchmark generator for combinatorial interaction testing tools.
\newblock \emph{Journal of Systems and Software}, 209:111920, 2024.
\newblock \url{https://doi.org/10.1016/j.jss.2023.111920}.

\bibitem{ctcompetition2023}
CT Competition organizers.
\newblock Second edition of the Combinatorial Interaction Testing Competition, 2023.
\newblock \url{https://fmselab.github.io/ct-competition/results/2023/index2023.html}.

\bibitem{citbenchmark2023}
FMSE Lab.
\newblock CIT Benchmark Generator: \texttt{INDUSTRIAL\_8} evaluation model, 2023.
\newblock \url{https://github.com/fmselab/CIT_Benchmark_Generator}.

\end{thebibliography}

\end{document}
